\section{Introduction}
\label{sec:intro}

Enterprise documents preserve regulations, failure analyses and operating
procedures. To answer a question from these materials, a retrieval system
must reach the right collection, identify a document and expose the passage
that supports the answer. Each step depends on earlier management decisions.
A collection description guides the search, a document must be listed on
the selected shelf, and the index must point to its current location.
An accurate passage ranker has little opportunity to recover evidence that
the access path never reaches.

Document acquisition, organisation and retrieval are often evaluated at
different boundaries. Enterprise content management studies document
governance \cite{paivarinta-ecm-2005,tyrvaainen-ecm-2006}, while metadata
quality and tagging determine how users discover those documents
\cite{park-metadata-quality-2009,golder-tagging-2006}. Distributed
information retrieval addresses the selection of collections
\cite{callan1995searching,si2003redde}. Retrieval-augmented generation
(RAG) and agentic retrieval add passage ranking and tool-guided evidence
access \cite{rag-lewis,dpr,modular-rag,mcp}.
HMRS connects these responsibilities. It makes collection membership,
document metadata and source references available from ingestion through
query time, so that an operator can trace a failed search to the stages
that supplied its candidates.

The Hierarchical Management Retrieval System (HMRS) is a self-hosted
knowledge document platform. An ingest agent uses parsed content to select
a collection and describe the document. Management operations update the
files, tree registry and library metadata, and reconciliation checks them
against one another. At query time, a catalogue-first search uses collection
and document descriptions to reach source passages. A pinned search begins
within a selected collection, and an optional server-side evidence gate
returns scored passages with source coordinates. The access policies use the
same managed state, so their choices can be traced independently of parsing
or index maintenance.

\noindent\textbf{Contributions.}
\begin{enumerate}
\item \textbf{A management model for the state behind retrieval.} HMRS
organises knowledge documents as collections, documents and segments, and
connects these managed objects to catalogue, shelf, document and segment
views. Collection membership, descriptions and source references link the
management of documents to the choices made during evidence access
(Section~\ref{sec:system}).
\item \textbf{A content-informed document lifecycle.} The acquisition
workflow uses parsed content to select a collection and derive document
descriptions and tags. Management operations and reconciliation checks
connect the document files, tree registry and per-library metadata, while
index and retrieval probes help expose stale state after changes
(Sections~\ref{sec:acquisition}--\ref{sec:organization}).
\item \textbf{Agent access across management levels.} MCP interfaces expose
the catalogue, document descriptions, source documents and indexed passages
to an agent. Catalogue-first and pinned policies let the agent narrow a
search using managed descriptions or start within selected collections
(Section~\ref{sec:retrieval}).
\item \textbf{A source-traceable evidence gate.} An optional server-side
gate fetches candidate text, scores its passages against the question, and
returns admitted passages with source references and error information.
Candidate ordering is a separate decision made by the calling workflow or
hybrid tool. Answer instructions distinguish a completed search without
support from an incomplete execution
(Sections~\ref{sec:gate} and~\ref{sec:lanes}).
\end{enumerate}

\subsection{Research questions and reading guide}
\label{sec:rq}

\textbf{RQ1} examines document acquisition, collection membership and
lifecycle consistency. \textbf{RQ2} measures collection selection and
document visibility under catalogue-first access. \textbf{RQ3} compares
target-string containment and retrieval cost across the tested evidence
configurations. \textbf{RQ4} compares the order of retained candidates and
examines the admission scores available in the archive.
\textbf{RQ5} examines abstention, runtime, tool cost and execution failures
at the answer layer.

Section~\ref{sec:related} positions HMRS within related work.
Sections~\ref{sec:system}--\ref{sec:implementation} describe the management
model and its implementation. Section~\ref{sec:eval} reports the evaluation,
followed by the discussion, limitations and conclusions in
Sections~\ref{sec:discussion}--\ref{sec:conclusion}.
